\section{Finiteness and Growth}

The cusp is never the complete infinite $\{4,3,4\}$ honeycomb. A real mesh holds finitely many cells and grows
one frontier layer at a time. So a fair question is how large the finite, growing cusp has to be before it looks
like today's space. The answer turns out to be surprisingly small, and the observed universe is vastly beyond it.
This section establishes that a finite chunk far short of infinity is more than enough, and it keeps the question
of whether the crystal is ultimately finite or infinite genuinely open.

\subsection{The complete cusp is never reached, and need not be}

The complete $\{4,3,4\}$ cusp is the infinite limit. A finite mesh never realizes it and never gets there in any
finite number of beats. This is not a flaw in the model. Physical space is never meant to be the complete cusp.
It is an ever-growing finite chunk of $\{4,3,4\}$, in exactly the way a finite crystal in a laboratory is
described by an infinite-lattice idealization while remaining finite. Local physics only sees the local lattice,
and the idealization is a convenience for reasoning about the bulk behaviour of a finite object.

The never-reaching has a reading that fits the rest of the theory. It is the eternal expansion and the open
future. The frontier always has more to add, and arriving at the complete cusp would be the end of time rather
than its fulfilment. The growth is the point. A finished crystal would be a finished universe, and the model does
not describe one.

\subsection{Convergence is fast}

To measure sufficiency we build cubic chunks of growing side $L$ and watch the observables settle. This is cheap
to do. The cusp's intrinsic structure is the $\mathbb{Z}^3$ cubic lattice, so a chunk of side $L$ is an
$O(L^3)$ object built directly, with no expensive hyperbolic horoball construction needed. We track two numbers,
the spectral dimension and the effective gravity exponent, as $L$ increases.

\begin{table}[h]
\centering
\begin{tabular}{rcc}
\toprule
side $L$ & spectral dimension & gravity exponent \\
\midrule
$5$  & $2.58$ & $2.06$ \\
$7$  & $2.88$ & $1.61$ \\
$11$ & $2.99$ & $1.34$ \\
$17$ & $2.99$ & $1.40$ (converged) \\
$25$ & $2.99$ & $1.42$ \\
$35$ & $2.99$ & $1.43$ \\
\bottomrule
\end{tabular}
\end{table}

The spectral dimension reaches three by about $L \sim 11$ cells. Local physics becomes continuum-like and
$L$-independent by about $L \sim 17$ cells \cite{pollard2026vibetest}. That is roughly $5000$ cells, and mapping
size to growth through a scale factor $a(t) \sim e^{Ht}$, it amounts to only three or four beats of growth. A
clean local physics is reached almost at once. The observed universe, by contrast, spans on the order of
$10^{61}$ cells, which is astronomically deep inside the continuum regime, so finite-size and edge effects at any
observable scale are utterly negligible. The figure of $10^{61}$ is reached by analytic extrapolation of the
$\mathbb{Z}^3$ lattice toward isotropic three-dimensional elasticity, as a power of $1/L$, rather than by
building that many cells.

\begin{result}[A small chunk suffices]
Cubic chunks of side $L$ reach spectral dimension $3$ by $L \sim 11$ and $L$-independent continuum physics by
$L \sim 17$, about $5000$ cells or a few beats of growth. The observed universe at $\sim 10^{61}$ cells is far
inside this regime, so finite-size effects at observable scales are negligible \cite{pollard2026vibetest}.
\end{result}

\subsection{The refining universe}

The growth has a shape that matches what cosmology sees. The universe expanded fast to a rough geometry, a small
chunk with measurable departures from perfect three-dimensionality, and it refined toward clean three-dimensional
space as it grew. The large-scale flatness and isotropy improve with expansion, approaching but never reaching
the perfect infinite-cusp limit. This is consistent with the observed picture of a universe flat to about
$0.4$ percent and of finite age. The early roughness and the late smoothness are two readings of the same
convergence table, taken at small and large $L$.

One quantity does not vanish with growth. The fixed cubic anisotropy, the order-four imperfection of the cusp's
octahedral point group, is a permanent feature rather than a finite-size artifact. It does not refine away as the
chunk enlarges, because it is built into the lattice symmetry itself and not into its size. The later audit
treats this residual directly. Here the point is only that the refining behaviour and the permanent anisotropy
are distinct, and the convergence table separates them.


A finite cusp, far from infinity, is more than sufficient to look like today's space. The continuum regime
arrives in a few beats, and the real universe is enormously beyond the threshold. Whether the crystal is ultimately
finite or infinite is left open, and nothing in the convergence result depends on settling it, since either way
the observable patch sits deep inside the continuum.
